% !TEX root = ../main.tex

% math typesetting utilities
\newcommand*{\given}{\;\middle|\;}
\newcommand*{\hateq}{\mathrel{\widehat{=}}}
\newcommand*{\tuple}[1]{\left\langle #1 \right\rangle}
\newcommand*{\set}[1]{\left\{ #1 \right\}}
\newcommand*{\multiset}[1]{\left\{\!\left\{ #1 \right\}\!\right\}}
\newcommand*{\evalat}[2][]{\left. #2 \right\rvert_{#1}}
\newcommand*{\bmat}[1]{\begin{bmatrix} #1 \end{bmatrix}}
\newcommand*{\sbmat}[1]{\begin{bsmallmatrix} #1 \end{bsmallmatrix}}
\newcommand*{\dbmat}[1]{\everymath{\displaystyle} \begin{bmatrix} #1 \end{bmatrix}}
\newcommand*{\abs}[1]{\left\lvert #1 \right\rvert}
\newcommand*{\norm}[1]{\left\lVert #1 \right\rVert}
\newcommand*{\bvec}[1]{{\symbfit{#1}}}

% math operators
\DeclareMathOperator*{\argmax}{arg\,max}
\DeclareMathOperator*{\argmin}{arg\,min}
\DeclareMathOperator*{\Exp}{\symbb{E}}
\DeclareMathOperator*{\Prob}{\symbb{P}}
\newcommand*{\expectation}[2][]{\Exp_{#1} \left[ #2 \right]}
\newcommand*{\probability}[1]{\Prob \left[ #1 \right]}

% environments
\newtheorem{theorem}{Theorem}
\newtheorem{problem}{Problem} \crefname{problem}{Problem}{Problems}
\newtheorem{hypothesis}{Hypothesis} \crefname{hypothesis}{Hypothesis}{Hypotheses}
